\documentclass[letterpaper, 10 pt, conference]{ieeeconf}
\IEEEoverridecommandlockouts
\overrideIEEEmargins

\usepackage{graphicx}
\usepackage[T1]{fontenc}   % needed for accents such as \k{a} in bib entries
% ieeeconf.cls defines \proof; amsthm redefines it -> "already defined".
\let\proof\relax
\let\endproof\relax
\usepackage{amsmath, amssymb, amsfonts, amsthm}
\usepackage[dvipsnames]{xcolor}
\usepackage{array, tabularx, booktabs, multirow, makecell}
\usepackage[ruled,linesnumbered,lined]{algorithm2e}
% Captions as in IEEEtran+caption: "Fig. 1. text" / "TABLE I. text", tables captioned below.
\usepackage[font=footnotesize,labelsep=period,skip=5pt]{caption}
\usepackage{cleveref}

\graphicspath{{figures/}{figures/final/}{figures/}}
% The figures in figures/finalfinal/ are hand-finished exports carrying ~7pt text on a
% ~2.5in canvas. Stretched to a full column their text renders near 10pt, larger than the
% 8pt captions, so every figure goes in through one of these two lengths instead.
\newlength{\figw}\setlength{\figw}{0.80\columnwidth}      % single-column figures
\newlength{\figwide}\setlength{\figwide}{0.75\textwidth}  % Fig. 2, the two-column float

\newtheorem{theorem}{Theorem}
\newtheorem{proposition}{Proposition}
\newtheorem{lemma}{Lemma}
\theoremstyle{remark}
\newtheorem*{note}{Note}

\crefname{figure}{Fig.}{Fig.}
\Crefname{figure}{Figure}{Figures}
\Crefname{equation}{Equation}{Equations}
\crefname{theorem}{Theorem}{Theorems}
\crefname{section}{Section}{Sections}
\crefname{proposition}{Proposition}{Propositions}
\crefname{algorithm}{Algorithm}{Algorithms}
\crefname{algocf}{Algorithm}{Algorithms}
\crefname{table}{Table}{Tables}


\definecolor{MyPurple}{HTML}{800080}
\definecolor{MyOrange}{HTML}{d38100}
\definecolor{MyGreen}{HTML}{00bf20}
\definecolor{MyBlue}{HTML}{1455a6}

% ---------------------------------------------------------------------------
% \todo{...} — inline red note for drafting. Swap the two definitions below to
% hide every note at once for a submission build.
% ---------------------------------------------------------------------------
\newcommand{\todo}[1]{\textcolor{red}{\bfseries [TODO: #1]}}
%\newcommand{\todo}[1]{}

% ---------------------------------------------------------------------------
% DRAFT SCAFFOLDING — space budget. Delete this whole block, every \lipsum,
% every \figbox and every \budget before submission.
%   \budget{...}            blue marker: target length for the section it heads
%   \lipsum[a-b]            filler prose, ~100 words per paragraph
%   \figbox{<h>}{<label>}   grey box the height the final figure will occupy
%   \nocite{*}              pulls all of refs.bib in so the bibliography has
%                           real height (final list will be ~2x this long)
% Calibrated so the whole skeleton lands on exactly 8 pages. As you replace a
% \lipsum with prose, the page count tells you whether you are over or under.
% ---------------------------------------------------------------------------
\usepackage{lipsum}
\newcommand{\budget}[1]{\textcolor{RoyalBlue}{\bfseries [BUDGET: #1]}}
\newcommand{\figbox}[2]{%
  \fcolorbox{gray}{gray!12}{\parbox[c][#1][c]{\dimexpr\linewidth-2\fboxsep-2\fboxrule\relax}%
    {\centering\ttfamily\footnotesize #2}}}

\title{\LARGE \bf CONSORT: Cooperative Navigation with Support-Optimized Routing and Trajectories for GPS-Denied Escort Missions}

% ICRA 2027 uses double-anonymous review: the submitted PDF must carry no
% author names or affiliations. Restore the real list after acceptance.
\author{Tony Zhang$^{1*}$, Ryan Bartholomew$^{2*}$, Kalliyan Velasco$^2$, Joshua Mangelson$^2$, Brady Moon$^2$}

\begin{document}

\maketitle
\thispagestyle{empty}
\pagestyle{empty}

\begin{abstract}
Centralized multi-agent path planning with cooperative localization in GPS-denied environments is generally computationally 
intractable when formulated in full belief space due to tightly coupled uncertainty dynamics. Prior work either approximates belief dynamics
or restricts inter-agent coupling. We study a multi-agent planning problem in which a primary agent must reach a goal in minimum time 
subject to a chance-constrained collision-avoidance requirement and a terminal bound on localization uncertainty. We obtain a tractable surrogate
of the underlying belief-space problem by optimizing over the mean-state trajectories of all agents and propagating covariance for terminal constraint evaluation only. 
We propose CONSORT (Cooperative Navigation with Support-Optimized Routing and Trajectories), a discrete-to-continuous pipeline combining hexagonal-grid joint search with spline-based trajectory refinement. 
Across 50 randomized scenarios, when the terminal uncertainty constraint is 50\% of the terminal uncertainty of a reference lone primary agent plan, CONSORT
finds solutions in 80\% of the scenarios against 54\% for the strongest baseline. We validate these results through a high-fidelity marine robotics simulator and show that
the generated paths satisfy the collision constraint in 100\% and the terminal uncertainty bound in 93\% of closed-loop Monte Carlo trials.

\end{abstract}

% ===========================================================================
\section{INTRODUCTION}
% ===========================================================================
%\budget{$\approx$550 words, 0.7 page incl.\ Fig.~1}


In environments where GPS signal is unreliable, such as underwater, or in regions subject to spoofing or jamming, agents must rely on other forms of information
for localization. In these environments, autonomous agents struggle with state uncertainty due to inertial measurement unit (IMU)
drift. As the estimation error accumulates, an agent's pose uncertainty can grow to a point where waypoint tracking degrades and mission objectives are compromised.
Common planning strategies ameliorate this by biasing trajectories toward information-rich regions~\cite{roy_coastal_1999} and introducing cooperating agents to utilize inter-agent measurements for 
cooperative localization~\cite{roumeliotis2002dist,bahr2009coop}. Consequently, trajectory planning must account for uncertainty propagation when estimation error impedes task performance. Planning under uncertainty is naturally posed as a POMDP~\cite{kaelbling1998pomdp}, whose exact solution is generally intractable~\cite{papadimitriou1987mdp}. Gaussian belief-space planners recover tractability for a single agent by propagating beliefs along candidate trajectories~\cite{prentice2009brm, vandenberg2011lqgmp}, typically assuming the most likely observation is received~\cite{platt2010belief}. Extending this to a team exponentially grows the joint belief search space, and the beliefs cannot be factored as the inter-agent measurements couple them. This motivates surrogate formulations that optimize mean trajectories while propagating uncertainty only for constraint evaluation.

In this work, we consider a centralized multi-agent planning with cooperative localization (MAPP-CL) problem, a variant of the CL-MRMP problem of Theurkauf et al.~\cite{theurkauf2025coop}, in which a primary agent must reach a goal in minimum time while subjected to a chance-constrained collision-avoidance 
requirement and a terminal bound on localization uncertainty. We call this the escort problem: the support agents have no goal locations and are routed solely to hold the primary's terminal uncertainty under the constraint.

We propose CONSORT (Cooperative Navigation with Support Optimized Routing and Trajectories), which first routes the team over a joint discrete state space that captures motion constraints and inter-agent ranging, then lifts that route into continuous space as B-spline trajectories to be smooth and dynamically feasible.

\begin{figure}[t]
  \centering
  \includegraphics[width=\columnwidth]{figures/fig1 copy2.pdf}
  \caption{The escort problem. The primary must reach the goal with its
    localization uncertainty under a terminal bound, and the support is routed
    to hold it there. \textcolor{MyPurple}{Support agent} \textcolor{MyOrange}{localizes with landmark} and \textcolor{MyGreen}{relays~position~data}, enabling the \textcolor{MyBlue}{primary agent} to take a more direct path. Both trajectories are searched on the heading-aware
    hexagonal lattice and then refined as B-splines.}
  \label{fig:overview}
\end{figure}

Our contributions are as follows:
\begin{enumerate}
  \item A formulation of centralized MAPP-CL in the presence of obstacles, in which support
    agents are routed jointly with a primary agent subject to a terminal
    localization bound and chance-constrained collision avoidance under both
    agent and obstacle uncertainty.
  \item CONSORT, a discrete-to-continuous planner pairing a joint routing search over a heading-aware
    hexagonal lattice with B-spline refinement, that enforces the
    curvature, collision, and terminal-uncertainty constraints in both stages.
  \item A benchmark of 50 randomized scenarios against four baselines adapted
    from adjacent literature, on which CONSORT solves 80\% of scenarios at
    the 50\% constraint level against 54\% for the strongest baseline.
  \item Closed-loop Monte Carlo validation in HoloOcean~\cite{potokar2022holoocean} with an independent
    factor-graph estimator, showing that CONSORT's plans satisfy the collision constraint in 100\% of trials and the terminal
    uncertainty bound in 93\%.
\end{enumerate}

In the following sections, we formalize the problem, describe the discrete and continuous planning stages and evaluate our approach in a simulated GPS-denied environment.


% ===========================================================================
\section{RELATED WORK}
\label{sec:related}
% ===========================================================================
Planning under uncertainty is naturally posed as a partially observable Markov
decision process \cite{kaelbling1998pomdp}, which is intractable in general
\cite{papadimitriou1987mdp}.
Belief-space planners restrict beliefs to Gaussians and propagate them along
candidate trajectories, over a roadmap \cite{prentice2009brm}, a tree
\cite{bry2011rrbt}, or by ranking candidate paths under a closed-loop model \cite{vandenberg2011lqgmp},
usually under the assumption that the most likely observation is received
\cite{platt2010belief}. These methods plan for a single agent, and the
multi-robot case of LQG-MP plans the robots one at a time without
inter-agent measurements. With several
agents the joint belief and the search space grow exponentially with the
size of the team.

Cooperative localization lets agents share pose estimates and measure one
another to reduce dead-reckoning drift \cite{roumeliotis2002dist,bahr2009coop}
and is established in underwater navigation \cite{paull2014auv}. Because cooperating agents share
measurement history, naive fusion double counts information, and covariance
intersection \cite{julier1997cov} is the usual remedy \cite{carrillo2013ci}.
Planning with cooperative localization in the loop is more recent. Theurkauf
et al. \cite{theurkauf2025coop} grow a sampling-based tree over the joint
belief of a team, and Liu et al. \cite{liu2026risk} plan risk-aware
trajectories for underwater teams. Our support agents have no goals of their
own and are routed only to meet the primary's terminal bound, and our joint
search is exhaustive on a lattice rather than sampled.

Multi-agent path finding solves the discrete coordination problem exactly
with conflict-based search \cite{sharon2015cbs} or approximately by planning
the agents one at a time in priority order
\cite{erdmann1987multiple,silver2005cooperative}, but it treats positions as
known. Hexagonal grids give six equidistant neighbours and paths that are
easier to smooth \cite{duszak2021hexgrid}, and B-splines give smooth
trajectories whose curvature can be bounded \cite{feng2024bsrrt}. Obstacle
avoidance under position uncertainty extends the chance-constrained
formulation of Blackmore et al. \cite{blackmore2011chance} to uncertain
obstacle positions, and we bound each spline segment by its MINVO hull
\cite{tordesillas2022minvo} so that the constraint is linear in the control
points.

% ===========================================================================
\section{PROBLEM FORMULATION}
\label{sec:pf}
% ===========================================================================
Let $\mathcal{I} = \{1,2,\dots,P\}$ denote a team of $P$ agents, where agent $i=1$
is the primary agent tasked with reaching a goal location $G\in\mathbb{R}^2$.
All other agents $2,\dots,P$ are support agents whose sole task is to assist the
primary agent's localization. All agents move at a common constant speed $v$,
so minimizing time is equivalent to minimizing trajectory length, and we use
arc length $s = vt$ as a shared clock throughout. The environment contains $M$
static convex obstacles $\mathcal{O}_1,\dots,\mathcal{O}_M$, each defined as an
intersection of half-planes $a_j^{\top}x \le b_j$ with unit outward normals
$a_j$ and offsets $b_j$, and each located only up to a covariance $\Sigma_o$.
The environment also contains a set $\mathcal{L}$ of landmarks with known
positions $\ell\in\mathbb{R}^2$ and positional uncertainty $\Sigma_\ell$.

\subsection{Decision Variables}

We optimize over mean-state trajectories only. The decision variables are the
heading profile $\theta_i(s)$ of every agent and the primary's path length
$S_1$. Positions follow from the headings at the arc samples
$s_k = k\,\Delta s$ through
\begin{equation}
  x_i(s_{k+1}) = x_i(s_k) + \Delta s
  \begin{bmatrix} \cos\theta_i(s_k) \\ \sin\theta_i(s_k) \end{bmatrix}.
  \label{eq:dyn}
\end{equation}

\subsection{Covariance Propagation}

Let $\Sigma_i(s)$ denote the $2\times2$ position covariance of agent $i$.
Covariance is not a decision variable. It is propagated along the mean
trajectories and used only to evaluate constraints. Between measurements the
agents dead-reckon, and we model the drift as process noise whose variance
grows linearly with distance and is larger along the direction of travel than
across it:
\begin{equation}
  \Sigma_i^{-}(s_{k+1}) = \Sigma_i(s_k) + Q\big(\theta_i(s_k)\big)\,\Delta s,
  \label{eq:prop}
\end{equation}
where $Q(\theta) = R(\theta)\,\mathrm{diag}(q_\parallel, q_\perp)\,R(\theta)^{\!\top}$,
$R(\theta)$ is the planar rotation, and $q_\parallel \ge q_\perp$ are the
along-track and cross-track noise intensities.

\subsection{Landmark Measurements}

A landmark $\ell$ provides a range-bearing measurement whose quality decays
smoothly with the distance $d_{i\ell} = \lVert x_i - \ell \rVert$ through the
logistic taper
\begin{equation}
  q_{i\ell} = \frac{1}{1 + \exp\!\big((d_{i\ell} - r_v)/w_v\big)},
  \label{eq:taper}
\end{equation}
which is close to one inside the visibility range $r_v$ and rolls off over a
width $w_v$, with no hard cutoff. Every landmark contributes information
weighted by its quality, and the posterior follows the information form of
the Kalman update for a direct position measurement, identical to the Joseph
form for $H = I$:
\begin{equation}
  \Sigma_i(s_k) = \Big( \Sigma_i^{-}(s_k)^{-1}
  + \sum_{\ell \in \mathcal{L}} q_{i\ell}\,\big(S_{i\ell} + \Sigma_\ell\big)^{-1} \Big)^{-1},
  \label{eq:lm}
\end{equation}
where $S_{i\ell} = R(\beta_{i\ell})\,\mathrm{diag}(\sigma_r^2, \sigma_b^2)\,R(\beta_{i\ell})^{\!\top}$
is the sensor noise rotated to the line-of-sight bearing $\beta_{i\ell}$.

\subsection{Inter-Agent Fusion}

Agents exchange position estimates at fixed arc-length intervals $\Lambda$.
Their estimates share history, so adding the information directly would
double count it, and we fuse with covariance intersection
\cite{julier1997cov}, which is consistent under unknown cross-correlation.
When agents $i$ and $j$ communicate at separation $d_{ij}$, agent $i$
receives $P_{j \to i} = (\Sigma_j + R_c)/w_{ij}$, where $R_c = \sigma_c^2 I$
is the communication noise and $w_{ij}$ is a taper of the form
\eqref{eq:taper} with half weight at the communication range $r_c$, and
updates
\begin{equation}
  \Sigma_i^{+} = \Big( \omega\,\Sigma_i^{-1} + (1-\omega)\,P_{j \to i}^{-1} \Big)^{-1},
  \label{eq:ci}
\end{equation}
with $\omega \in [0,1]$ chosen to minimize $\det\Sigma_i^{+}$. Agent $j$ is
updated in the same way, and as $w_{ij} \to 0$ the transferred information
vanishes.

\subsection{Optimization Problem}

The planning problem collects the objective and all constraints as
\begin{subequations}
\label{eq:opt}
\begin{align}
  & \min_{\{\theta_i(\cdot)\},\, S_1}\;\; S_1 \nonumber \\
  & \mathrm{s.t.}\;\;
    \det\big(\Sigma_1(S_1)\big)^{1/4} \le \bar{u}, \label{eq:term} \\
  & a_j^{\top} x_i(s) \ge b_j + r_a
      + z_\delta \sqrt{a_j^{\top}\big(\Sigma_i(s) + \Sigma_o\big)a_j},
      \label{eq:cc} \\
  & |\kappa_i(s)| \le 1/r_{\min}, \label{eq:kin} \\
  & x_i(0) = x_i^0,\;\; \Sigma_i(0) = \Sigma^0,\;\; x_1(S_1) = G.
      \label{eq:bc}
\end{align}
\end{subequations}

\textbf{Terminal uncertainty \eqref{eq:term}.} The primary must arrive with
bounded localization uncertainty. $\det(\Sigma)^{1/4}$ equals the standard
deviation $\sigma$ of an isotropic covariance $\sigma^2 I$, so $\bar{u}$ is
in metres.

\textbf{Obstacle avoidance \eqref{eq:cc}.} Every agent must clear every
obstacle with probability at least $1-\delta$ at every instant, so for each
agent $i$, obstacle $o$, and arc position $s$ the constraint must hold for at
least one face $j$ of $\mathcal{O}_o$. The margin covers the vehicle radius
$r_a$ and the positional uncertainty of both agent and obstacle, with
$z_\delta = \Phi^{-1}(1-\delta)$ and $\Phi$ the standard normal distribution
function \cite{blackmore2011chance}. A union bound over the $PM$
agent-obstacle pairs bounds the instantaneous probability of any collision by
$\Delta = PM\delta$.

\textbf{Kinematics \eqref{eq:kin}.} Each trajectory respects the minimum
turn radius $r_{\min}$, where $\kappa_i = \mathrm{d}\theta_i/\mathrm{d}s$ is
the curvature.

\textbf{Boundary conditions \eqref{eq:bc}.} All agents start from known
positions with a common prior covariance, and the primary must end at the
goal.

The belief dynamics do not appear as constraints: $\Sigma_i(s)$ is a
deterministic function of the decision variables through
\eqref{eq:dyn}--\eqref{eq:ci}, left implicit in \eqref{eq:opt}. The agents
are coupled only through \eqref{eq:ci}, so the supports enter the problem
only by shaping the primary's covariance.

% ===========================================================================
\section{METHOD}
\label{sec:method}
% ===========================================================================
Our planner solves \eqref{eq:opt} in two stages, a joint A* search over a
heading-aware hexagonal lattice and a continuous B-spline refinement of the
resulting plan, summarized in Algorithm~\ref{alg:joint}.

\subsection{Heading-Aware Hexagonal Graph}

A corridor of hexagonal cells \cite{duszak2021hexgrid} of width
$w_{\mathrm{hex}}$ covers the line from start to goal. A graph node pairs a
cell with
one of six headings spaced $60^\circ$ apart, and three moves leave each
node: continue straight, or turn by $\pm 60^\circ$ and step into the cell
faced by the new heading. The heading changes by at most $60^\circ$ per
cell, which is the graph-level counterpart of \eqref{eq:kin}. Landmarks are
appended as sensor nodes without edges that enter the search only through
\eqref{eq:lm}, a goal node absorbs all six headings of the goal cell so that
the arrival heading is free, and obstacles act on beliefs during the search
rather than removing cells. All-pairs shortest-path distances are
precomputed on the graph and serve as the heuristic.

\subsection{Joint Discrete Search}

The first stage runs an A*-style search \cite{hart1968astar} over the joint
state of the team. A search state stores, for every agent, its graph node,
its covariance, and the arc distance travelled. Each expansion advances
every agent by one edge, so the successors of a state are the Cartesian
product of the agents' three moves. The accumulated cost $g$ is the
primary's arc distance alone, and states are expanded in order of
$f = g + (1+\varepsilon)\,h$, where $h$ is the precomputed shortest-path
distance from the primary's node to the goal. Supports contribute no
heuristic cost, so $h$ is admissible, and $\varepsilon > 0$ gives a weighted
search \cite{pohl1970weighted} that trades a factor of at most
$(1+\varepsilon)$ on the seed length \cite{pearl1984heuristics} for fewer
expansions. Ties are broken toward lower primary
uncertainty and then toward fewer idle supports.

During expansion each agent's covariance is propagated along its edge by
\eqref{eq:prop}--\eqref{eq:lm}, and \eqref{eq:ci} is applied whenever two
agents reach a communication checkpoint at the same arc distance. A
successor is discarded if any agent violates \eqref{eq:cc} along its edge or
if a support travels farther than the primary. Pareto pruning keeps the
joint space manageable: at each joint node signature the search keeps a
frontier of $(g, \Sigma_1)$ labels and discards a new label if an existing
one has no larger $g$ and a covariance no larger in the positive
semidefinite order. The pruning is sound because the heuristic is shared at
the same node tuple and a larger covariance cannot become smaller under
further propagation.

The search is anytime with respect to \eqref{eq:term}. When a popped state
places the primary at the goal and the spline seed it induces
(Sec.~\ref{sec:cont}) clears all obstacles, the state is returned if
\eqref{eq:term} holds and otherwise kept as the incumbent if its terminal
uncertainty is the lowest seen so far. If the expansion budget is exhausted,
the incumbent is handed to the continuous stage, which attempts to recover
feasibility. Only the terminal bound is relaxed in this way, never obstacle
clearance.

\begin{note}
Six headings quantize direction to within $30^\circ$, so a lattice path that
traces a smooth trajectory is at most $2/\sqrt{3}$ times longer than it. If
\eqref{eq:opt} has a solution with strictly positive constraint slacks, a
fine enough lattice therefore contains a feasible seed within that factor of
the optimum, and the admissible search returns one no longer than it. The
factor comes from the move set rather than the resolution, and removing it
is the job of the refinement.
\end{note}

\begin{algorithm}[t]
\SetInd{0.4em}{0.8em}
\DontPrintSemicolon
\small
\SetKwInOut{Input}{Input}\SetKwInOut{Output}{Output}
\Input{hex graph $\mathcal{G}$, start states $x_i^0$, $\Sigma^0$, goal $G$,
       bound $\bar{u}$, expansion budget $N$}
\Output{B-spline control points $x$ of every agent}
$h(\cdot) \leftarrow$ \textsc{AllPairsShortestPaths}$(\mathcal{G})$\;
$\mathit{open} \leftarrow \{\text{joint start}\}$;
  $\mathit{inc} \leftarrow \varnothing$; $\mathit{seed} \leftarrow \varnothing$\;
\While{$\mathit{open} \neq \varnothing$ and expansions $< N$}{
  pop the state minimizing $f = g + (1+\varepsilon)\,h$\;
  \If{primary at $G$ and \textnormal{\textsc{SeedSplineClear}}$(\mathrm{state})$}{
    \lIf{$\det(\Sigma_1)^{1/4} \le \bar{u}$}{$\mathit{seed} \leftarrow \mathrm{state}$; \textbf{break}}
    \lElseIf{$\det(\Sigma_1)^{1/4}$ lowest so far}{$\mathit{inc} \leftarrow \mathrm{state}$}
  }
  \ForEach{joint successor (every agent one edge)}{
    propagate $\Sigma_i$ by \eqref{eq:prop}--\eqref{eq:lm}; fuse by \eqref{eq:ci}\;
    discard if \eqref{eq:cc} fails or a support outruns the primary\;
    insert unless dominated in $(g, \Sigma_1)$ at its signature\;
  }
}
\lIf{$\mathit{seed} = \varnothing$}{$\mathit{seed} \leftarrow \mathit{inc}$, the least-violating candidate}
$x \leftarrow$ control points of the seed splines;
  $\mathit{feas} \leftarrow \min s(x) \ge 0$\;
\ForEach{barrier stage; Adam step on $\nabla_x J(x;\mu)$}{
  accept iff $J$ decreases and ($\lnot\mathit{feas}$ or trial feasible)\;
  $\mathit{feas} \leftarrow \mathit{feas} \lor \text{trial feasible}$\;
  track least-violating and shortest feasible incumbents\;
  decay $\mu$ only once $\mathit{feas}$\;
}
\Return the best incumbent trajectory\;
\caption{Discrete-to-continuous escort planning}
\label{alg:joint}
\end{algorithm}

% Fig.~2 is a two-column float; LaTeX places figure* at the top of the NEXT
% page after its source position, so it is issued here to land on the page
% where Results begins. Move the source if it drifts.
\begin{figure*}[t]
  \centering
  \includegraphics[width=\figwide]{fig2}
  \caption{Constraint ladder on one scenario: a 1700\,m corridor of staggered
    walls, 14 obstacles and three landmarks. The terminal bound tightens from
    (a) to (d), and every tightening is paid for in primary length; the levels
    and the solutions they produce are in Table~\ref{tab:ladder}. In (a) the
    support climbs into the north pocket for the weak landmark there and
    rejoins; in (b) it takes the southern lane to a stronger landmark and
    relays it from behind the last wall; in (c) the primary adds a turn before
    the goal so that the support can relay from closer; in (d) it dips to the
    centre line so that the support can reach the last landmark in time.}
  \label{fig:ladder}
\end{figure*}

\begin{table}[t]
  \centering
  \footnotesize
  \setlength{\tabcolsep}{5pt}
  \begin{tabular}{@{}lcccc@{}}
    \toprule
    Panel & \makecell{Constraint level\\(\% of $U_\mathrm{ref}$)}
          & \makecell{Terminal\\bound $\bar{u}$ (m)}
          & \makecell{Primary\\length (m)}
          & \makecell{Terminal\\$\sigma$ (m)} \\
    \midrule
    (a) & 100 & 13.75 & 1741 & 9.73 \\
    (b) &  70 &  9.62 & 1761 & 8.89 \\
    (c) &  50 &  6.87 & 1822 & 6.66 \\
    (d) &  30 &  4.12 & 1890 & 3.91 \\
    \bottomrule
  \end{tabular}
  \caption{The four panels of Fig.~\ref{fig:ladder}. The constraint level is
    the terminal bound as a fraction of the uncertainty an unconstrained
    single agent reaches on the same scenario, which is 13.75\,m over a
    reference path of 1712\,m.}
  \label{tab:ladder}
\end{table}

\subsection{Continuous B-Spline Refinement}
\label{sec:cont}

The seed is lifted into continuous space as clamped cubic B-splines
\cite{piegl1997nurbs}, one per agent, with the seed's cell centres as control points. Support paths are padded to the primary's node count, and
the primary's final control point is pinned to $G$. The decision vector $x$
stacks the free control points: the interior points of the primary and all
but the first point of each support, whose terminal position is free.

The constraints of \eqref{eq:opt} form a slack vector $s(x)$, feasible when
$s(x) \ge 0$: the terminal bound, evaluated by propagating
\eqref{eq:prop}--\eqref{eq:ci} along the sampled splines, the curvature
bound, sampled from the spline derivatives, the requirement that no support
spline be longer than the primary's, and obstacle clearance. For the latter,
each spline segment is enclosed in the convex hull of its MINVO control
points \cite{tordesillas2022minvo}, and each segment-obstacle pair is
committed to the face of \eqref{eq:cc} that the seed clears by the largest
margin, which fixes the homotopy class and replaces the disjunction over
faces with linear rows, one per hull vertex. The refinement then minimizes
the primary's length under the log-barrier merit function
\begin{equation}
  J(x;\mu) = S_1(x) + \sum\nolimits_k B\big(s_k(x);\mu\big),
  \label{eq:merit}
\end{equation}
\begin{equation}
  B(s;\mu) =
  \begin{cases}
    -\mu \log s, & s > \delta_b, \\[2pt]
    -\mu \big( \log\delta_b + \tfrac{s-\delta_b}{\delta_b} \big), & s \le \delta_b,
  \end{cases}
  \label{eq:barrier}
\end{equation}
where the tangent extension below the knee $\delta_b$ keeps $J$ finite for
violated constraints and supplies a constant restoring gradient, so one
objective serves feasible and infeasible iterates alike.

We minimize $J$ with Adam \cite{kingma2015adam} on forward-difference
gradients and a backtracking line search over a schedule of barrier stages.
While no feasible iterate has been accepted, as when the search handed over
its incumbent, any step that decreases $J$ is accepted and $\mu$ is held
large so that the barrier drives the violation down. The first feasible
iterate arms a latch: from then on only feasible steps that decrease $J$ are
accepted and $\mu$ decays between stages, shifting the effort from
constraint satisfaction to shortening. The least-violating and the shortest
feasible iterates are tracked throughout, and the better one is returned, so
the output is never longer than a feasible seed nor more violating than an
infeasible one.

% ===========================================================================
\section{RESULTS}
\label{sec:results}
% ===========================================================================

\subsection{Experimental Setup}

We evaluate the planner on randomized corridor scenarios with one primary and
one support. Each scenario places the goal 600 to 900\,m from the start,
draws two to four landmarks with map covariances of 0.6 to 1.35\,m standard
deviation, placed 280 to 430\,m to the side of the direct line on
alternating sides so that every fix costs a detour, and draws one to three
convex obstacles of up to 220\,m across, centred on a triangular
distribution so that most scenarios have an obstacle on the direct line. The
lattice width is 150\,m, the minimum turn radius 40\,m, the visibility range
100\,m with a 15\,m taper, the communication interval 100\,m, and the
communication range 300\,m with a 20\,m taper. Dead-reckoning variance grows
at $q_\parallel = 0.25$\,m$^2$ per metre along track, a standard deviation
of 5\,m after 100\,m, and at one ninth of that across track. The obstacle
risk is $\delta = 0.05$ per agent-obstacle pair with a 5\,m vehicle radius,
the search runs with $\varepsilon = 0$ and $10^6$ expansions, and the
refinement runs 15 barrier stages sharing 1000 Adam iterations.

For every scenario we first plan a single unconstrained agent, which gives a
reference length $L_\mathrm{ref}$ and terminal uncertainty $U_\mathrm{ref}$.
The bound is then tightened from 100\% to 30\% of $U_\mathrm{ref}$ in steps
of 10\%, and every planner is run at each level. A planner that fails at one
level is counted as failing at the tighter ones. Scenarios in which a
straight line already meets the loosest bound pose no localization problem
and are discarded, which removed 34 of 84 draws to leave the 50 scenarios
reported here. Lengths are reported relative to $L_\mathrm{ref}$ over the
scenarios a planner solved, and wall-clock times are for whole Julia
processes including start-up and graph construction.

\subsection{Baselines}

We compare against four planners adapted from adjacent literature. All share
the covariance model of Sec.~\ref{sec:pf}, the obstacle test, and the
refinement of Sec.~\ref{sec:cont}, so the comparison isolates how the team
is routed.

\textit{Greedy.} The primary commits to its shortest obstacle-clear route
without uncertainty reasoning, and at every step the support picks, among
its three successors, the move that most reduces the primary's uncertainty
at that step, the one-step information-gain rule of greedy exploration
\cite{bourgault2002info}. This removes lookahead from the support.

\textit{Formation.} The primary plans alone with a weighted A* over cells
and beliefs while the support holds a fixed slot in the primary's body
frame, one cell out at a $60^\circ$ multiple off its heading, as in
leader-referenced formation control \cite{balch1998formation}. Goal candidates
are accepted on the uncertainty the primary has with the support attached,
and both lateral slots are tried. This removes support routing.

\textit{Sequential.} Prioritized planning
\cite{erdmann1987multiple,silver2005cooperative}. The primary is planned
first with the support parked at its start pose, under a bound slackened to
twice the threshold, the support is then planned against the committed
primary to minimize the primary's terminal uncertainty, and the primary is
re-planned once against the finished support path. This removes the joint
state space.

\textit{CL-GBT.} Our cooperative-localization-guided belief tree, in the
spirit of Theurkauf et al. \cite{theurkauf2025coop}, who extend Gaussian
belief trees \cite{ho2022gbt} to grow one tree over the joint belief of a
team. We keep that structure, one rapidly exploring random tree over the
joint state in continuous space, with goal-biased sampling for the primary
and a cloning bias that places the support's sample on the primary's. It is
not a port. Their joint Kalman belief with cross-correlations, their bound
on the probability that a planned inter-agent measurement falls out of
range, their robot-robot collision constraint and their re-branching step
are replaced by our belief model, which credits every scheduled fusion at a
range-tapered weight. Its iteration budget is set so that its compute
ceiling matches ours. This replaces exhaustive lattice search with sampling.

\subsection{Comparison With Baselines}

\begin{figure}[t]
  \centering
  \includegraphics[width=\columnwidth]{fig3}
  \caption{Baseline comparison over the 50 randomized scenarios as the terminal
    uncertainty bound tightens from 100\% to 30\% of the single-agent
    reference $U_\mathrm{ref}$. Top: fraction of scenarios solved. Bottom:
    mean primary path length over the solved scenarios, normalized by the
    single-agent reference length $L_\mathrm{ref}$ (standard deviations in
    Table~\ref{tab:sweep}).}
  \label{fig:length_success}
\end{figure}

\begin{figure}[t]
  \centering
  \includegraphics[width=\columnwidth]{fig4_wall}
  \caption{Wall-clock time per planner over successful runs at five constraint
    levels (log scale). Boxes span the interquartile range, whiskers extend
    to $1.5\times$ IQR, and runs beyond the whiskers are omitted. Levels with
    fewer than five solved scenarios are drawn as individual runs. Only
    CL-GBT grows with the bound: at 60\% its upper quartile reaches 61\,s and
    its whisker 124\,s.}
  \label{fig:wall}
\end{figure}

Fig.~\ref{fig:length_success} reports success rate and path length against
the constraint level, and Table~\ref{tab:sweep} lists the lengths with their
spread. Every planner but Greedy solves nearly every scenario at the loosest
level, and the planners separate as the bound tightens. Ours solves every
scenario down to 70\%, 96\% at 60\%, 80\% at 50\%, and 40\% at 30\%.
Sequential matches it down to 80\% but falls to 52\% at 50\% and 16\% at
30\%, and its paths grow to more than twice the reference length, since at a
tight bound the primary must detour to landmarks on its own before the
support is planned. Formation degrades more gracefully and is the strongest
baseline from 50\% down, with 54\% at half the reference bound and 22\% at
30\%. CL-GBT solves half the scenarios at 60\% and 6\% at 30\%, with paths
1.5 to 2 times the reference length because a random tree returns the first
path it finds. Greedy solves 90\% at the loosest level and almost nothing
below it, since a support without lookahead rarely reaches a landmark in
time.

The price of the higher success rate is modest. Our paths are the shortest
among the solved scenarios at every level, from $1.14\,L_\mathrm{ref}$ at
100\% to $1.41\,L_\mathrm{ref}$ at 30\%, with the smallest spread across
scenarios. Fig.~\ref{fig:wall} reports wall-clock time: a median of 18\,s
per run against 14\,s for Formation and Sequential and 16\,s for the
single-agent reference, most of it shared overhead. Only CL-GBT grows with
the bound, and its longest run took 730\,s.

\subsection{Ablation}

\begin{figure}[t]
  \centering
  \includegraphics[width=\columnwidth]{fig5}
  \caption{Ablation on a separate set of 30 scenarios drawn without the
    straight-line screen. Same axes as Fig.~\ref{fig:length_success}: top,
    success rate; bottom, mean path length over the solved scenarios
    (standard deviations in Table~\ref{tab:sweep}).}
  \label{fig:abl_length_success}
\end{figure}

\begin{figure}[t]
  \centering
  \includegraphics[width=\columnwidth]{fig6}
  \caption{Wall-clock time of the ablation arms over successful runs, drawn as
    in Fig.~\ref{fig:wall} on a linear axis. The refinement adds about
    1.5\,s to the discrete stage. The straight seed's single run at 60\% is
    drawn as a tick.}
  \label{fig:abl_wall}
\end{figure}

\begin{table}[t]
      \centering
  \footnotesize
  \setlength{\tabcolsep}{4pt}
  \begin{tabular}{@{}lcccc@{}}
    \toprule
    & \multicolumn{4}{c}{Constraint level (\% of $U_\mathrm{ref}$)} \\
    \cmidrule(l){2-5}
    Method & 100 & 70 & 50 & 30 \\
    \midrule
    Ours          & 1.14$\pm$0.09 & 1.22$\pm$0.23 & 1.34$\pm$0.26 & 1.41$\pm$0.21 \\
    Formation     & 1.22$\pm$0.12 & 1.34$\pm$0.25 & 1.33$\pm$0.16 & 1.48$\pm$0.45 \\
    Sequential    & 1.18$\pm$0.14 & 1.41$\pm$0.92 & 1.36$\pm$0.50 & 2.76$\pm$1.51 \\
    CL-GBT        & 1.47$\pm$0.33 & 1.78$\pm$0.43 & 2.01$\pm$0.41 & 2.07$\pm$0.26 \\
    Greedy        & 1.24$\pm$0.10 & 1.21$\pm$0.06 & 1.13$\pm$0.01 & 1.13$\pm$0.01 \\
    \midrule
    Full pipeline          & 1.12$\pm$0.14 & 1.17$\pm$0.23 & 1.22$\pm$0.22 & 1.29$\pm$0.22 \\
    Discrete only          & 1.20$\pm$0.15 & 1.23$\pm$0.23 & 1.25$\pm$0.22 & 1.29$\pm$0.22 \\
    Straight seed + cont.\ & 0.99$\pm$0.04 & 1.00$\pm$0.00 & -- & -- \\
    \bottomrule
  \end{tabular}
  \caption{Mean $\pm$ SD of the primary path length over the solved scenarios,
    normalized by $L_\mathrm{ref}$, at four constraint levels. Top: the 50
    screened scenarios of Figs.~\ref{fig:length_success}--\ref{fig:wall}.
    Bottom: the 30-scenario ablation set of
    Figs.~\ref{fig:abl_length_success}--\ref{fig:abl_wall}. ``--'': no
    scenario solved.}
  \label{tab:sweep}
\end{table}

The ablation asks what each stage contributes. It runs on a separate set of
30 scenarios drawn without the straight-line screen, since one arm is the
straight seed itself, and compares the full pipeline with the discrete stage
alone and with the refinement seeded from a straight line
(Fig.~\ref{fig:abl_length_success} and the bottom block of
Table~\ref{tab:sweep}). The discrete stage alone solves the same scenarios
down to 70\% and one or two fewer at 60\% and 50\%, with paths 2 to 7\%
longer at every level down to 50\%, which is the length the refinement
recovers by cutting the lattice corners. The straight seed solves 53\% of
the scenarios at the loosest level and none below 60\%, and every solution
it returns has a length ratio of 1.00: the refinement never found a detour
from a straight seed and succeeded only where the straight line was already
feasible. The joint search finds the landmarks, and the
refinement makes the route short, at a cost of about 1.5\,s
(Fig.~\ref{fig:abl_wall}).

% ===========================================================================
\section{SIMULATION VALIDATION}
\label{sec:sim}
% ===========================================================================
Every planner in Sec.~\ref{sec:results} propagates the same covariance
model, so the benchmark cannot tell whether that model is calibrated. We
therefore fly a planned trajectory with a vehicle model, a controller, and a
state estimator that share none of the planner's equations, and check
whether the certified constraints hold over repeated runs with different
noise seeds.

\subsection{Simulator and Vehicle}

We use HoloOcean \cite{potokar2022holoocean,potokar2024holoocean}. Each
agent is a torpedo-shaped autonomous underwater vehicle with
six-degree-of-freedom Fossen dynamics \cite{fossen2021handbook}, driven by
the depth and heading autopilot of HoloOcean~2.0 \cite{romrell2025holoocean2}
at 30\,Hz at about 2.2\,m/s. An outer-loop waypoint follower
tracks the planned B-spline, resampled every 25\,m, with a 15\,m arrival
radius set by the turn radius, and it steers on the estimated position,
dead-reckoned forward between estimator updates, so estimation error feeds
back into the executed path. The planar field is placed on one depth plane
with the agents staggered in depth, so only agent-obstacle collisions are
under test. Three cadences are scheduled on arc length accumulated from
ground truth: the controller runs every tick, the estimator adds a pose node every 5\,m, and
an acoustic communication is attempted every 100\,m and succeeds only within
the planner's communication range. Measurements are ground truth with
additive Gaussian noise, so the simulation tests the consistency of the
belief model rather than the realism of any sensor.

\subsection{Factor-Graph State Estimator}
% Drafted from multiagent_base (factor_graph.py, main.py) at the state the
% shipped Monte Carlo set was flown, and awaiting review by the author of that
% stack. Credit for it belongs in the acknowledgments once the submission is
% de-anonymized.

The estimator is an existing GTSAM factor-graph stack
\cite{kaess2012isam2,dellaert2017factor}, developed independently of the
planner; what we supply is its configuration and, through
\cref{tab:noise}, its noise. It is an incremental smoother
\cite{kaess2012isam2} over $SE(3)$ pose nodes, one per 5\,m of arc for each
agent, returning the trajectory $\hat{\mathcal X}$ that minimizes the weighted
residuals $r_i$ of the factors $\mathcal F$ the graph holds,
\begin{equation}
  \hat{\mathcal X}=\arg\min_{\mathcal X}\;\sum_{i\in\mathcal F}
    \big\lVert r_i(\mathcal X)\big\rVert^{2}_{\Sigma_i}.
  \label{eq:fg}
\end{equation}

The graph is built in two layers. Each agent carries a dead-reckoning graph of
its own, holding a between-factor for the odometry of each node, a unary depth
factor, a unary three-axis orientation factor, and a range and bearing factor
to every landmark inside the visibility range. Landmarks are not free
variables: each is pinned by a prior carrying the planner's map covariance, so
the map is known and the problem is localization rather than mapping. This
graph is what the agent steers on, and for a support it is also the source of
the pose and the covariance its acoustic packets carry. The primary holds a
second graph, and the covariance compared against the planner's is read from
that one. It takes the primary's own nodes as they are created, and a
support's nodes in a batch when a communication succeeds, seeded with the pose
that support solved for itself, together with one inter-agent factor per
communication. Measurements enter both layers once rather than being drawn
again, and every one of them is ground truth plus additive Gaussian noise:
odometry is synthesized from the simulator's pose sensor, and there is no
inertial or Doppler sensor model.

Three configuration choices decide what the comparison in
Sec.~\ref{sec:sim}-C measures, and all three were in force for the runs
reported there. Poses are estimated directly in the world frame: the start
pose is known from the plan, so no separate map-origin node is estimated, and
the covariance compared against the planner's is the $xy$ block of the
world-pose marginal. A successful communication delivers the sender's
buffered nodes into the primary's graph together with a relative-pose factor
between the two vehicles, whose translation noise carries the planner's own
link weight; a bearing-only factor would transfer almost no position
information at these separations, and the primary could never inherit the
support's landmark fix. Finally, every factor's noise is derived from the
planner's model rather than set independently, so that the two codebases
measure uncertainty with one ruler.

\begin{table}[t]
      \centering
  \scriptsize
  \setlength{\tabcolsep}{3pt}
  \begin{tabular}{@{}llc@{}}
    \toprule
    Factor & Variables & $\sigma$ \\
    \midrule
    Odometry (between, per 5\,m node) & $\mathbf T_{k-1},\mathbf T_k$ & 1$^\circ$/axis; 0.22, 0.02, 0.01\,m \\
    Depth (unary) & $\mathbf T_k$ & 0.5\,m \\
    Compass (unary) & $\mathbf T_k$ & 0.86$^\circ$/axis \\
    Landmark az./el./range & $\mathbf T_k,\ell$ & 1.05$^\circ$, 1.05$^\circ$, 2.09\,m; $r<100$\,m \\
    Inter-agent relative pose & $\mathbf T^{s}_k,\mathbf T^{p}_k$ & 0.86$^\circ$/axis; $0.5/\sqrt{w}$\,m \\
    Initial pose prior & $\mathbf T_0$ & 10$^\circ$/axis; 1.15, 0.76, 0.5\,m \\
    \bottomrule
  \end{tabular}
  \caption{Factor noise of the simulated estimator, derived from the planner's
    noise model so that both codebases share one ruler: along-track odometry
    from $q_\parallel$, the compass from $q_\perp$, landmark noise from the
    planner's per-pass value rescaled to one sighting per 5\,m node, and
    communication noise $\sigma_c$ divided by $\sqrt{w}$ for link weight $w$.
    Poses are estimated in the world frame; the prior on the known start pose
    carries the planner's $\Sigma^0$ in $x$, $y$.
    Depth is the one channel the estimator treats more conservatively than
    the simulation: its factor keeps the stack's own 0.5\,m against the
    0.01\,m actually injected.}
  \label{tab:noise}
\end{table}
\subsection{Monte Carlo Results}

\begin{figure}[t]
  \centering
  \includegraphics[width=\columnwidth]{fig7_holoocean_screenshot}
  \caption{The two vehicles in HoloOcean~\cite{potokar2022holoocean}. Each is a
    torpedo-shaped autonomous underwater vehicle flying its planned spline at a
    fixed depth over the seabed of the simulator's open-water world.}
  \label{fig:holoocean}
\end{figure}

\begin{figure}[t]
  \centering
  \includegraphics[width=\figw]{fig8}
  \caption{Closed-loop Monte Carlo of the behind-wall plan in HoloOcean, 30 runs
    with different noise seeds: planned trajectories (dashed) against the
    mean executed trajectory of each agent, with the $\pm2\sigma$
    cross-track envelope. The primary passes north of the wall while the
    support diverts south to the landmark and rejoins before the goal.}
  \label{fig:mc_paths}
\end{figure}

\begin{table}[t]
      \centering
  \footnotesize
  \setlength{\tabcolsep}{4pt}
  \begin{tabular}{@{}lcccc@{}}
    \toprule
    Case & Runs & \makecell{Collision\\rate ($\le\Delta$)} & \makecell{Terminal\\bound met} & \makecell{Pred.\,/\,emp.\\$\sigma_T$} \\
    \midrule
    Behind wall & 30 & 0/30 & 28/30 & 1.80\,/\,1.70 \\
    \bottomrule
  \end{tabular}
  \caption{Monte Carlo constraint satisfaction in HoloOcean for the behind-wall
    plan of Fig.~\ref{fig:mc_paths} (two agents, one obstacle, one landmark,
    $\bar u = 1.8$\,m), 30 runs. Collision rate counts runs with any obstacle
    penetration against the allocated $\Delta = 0.10$; the terminal bound is
    met when the estimator's marginal at the goal is under $\bar u$;
    $\sigma_T$ is the planner's prediction against the sample covariance of
    the true terminal error.}
  \label{tab:mc}
\end{table}

We flew the behind-wall plan of Fig.~\ref{fig:mc_paths} 30 times with
different noise seeds. The scenario hides its one landmark behind a wall
400\,m long, so only the support can see it, under a terminal bound of
1.8\,m against which the planner predicts 1.80\,m. Every run reached the
goal within 2\% of the
planned length and none penetrated the wall, against an allocated collision
probability of $\Delta = 0.10$ (Table~\ref{tab:mc}).
The estimator's terminal marginal was under the bound in 28 of 30 runs, with
the two exceptions at 1.85 and 1.96\,m. The sample covariance of the true
terminal error gives 1.70\,m, within 6\% of the prediction, while the
estimator's mean marginal of 1.51\,m is more optimistic than both. The true
error also carried a mean offset of 2.6\,m along track and 0.8\,m across, a
bias of the estimate that a covariance comparison does not capture. The
$2\sigma$ cross-track envelope in Fig.~\ref{fig:mc_paths} has a half-width
of about 4\,m along most of the 1.1\,km route and 8\,m on the primary's
final bend, with a median cross-track error of 2.5\,m RMS. The baselines are
not re-flown, because they share the belief model whose calibration is the
quantity under test.

% ===========================================================================
\section{DISCUSSION AND LIMITATIONS}
% ===========================================================================
The discrete stage is complete only with respect to the lattice. The lattice
is laid over the corridor without regard to the obstacles, so the lattice
paths that clear them are neither a subset nor a superset of the continuous
feasible set, and a route that exists only between two cells is not found.
The refinement can shorten a seed and recover a small violation but cannot
change its homotopy class. A finer lattice narrows the gap at the cost of a
larger graph, and the all-pairs heuristic is cubic in the number of nodes.

The belief model is Gaussian and assumes that the planned trajectory
receives the maximum-likelihood observation at each step
\cite{platt2010belief}. The simulation
tests calibration against an independent estimator fed ground-truth
measurements plus Gaussian noise, so Table~\ref{tab:mc} says the planner's
covariance is consistent with that estimator under the noise of
Table~\ref{tab:noise}, not that either models an acoustic modem or a sonar. The along-track bias in the terminal error shows that a consistent
covariance does not rule out a biased estimate.

The plan is centralized, computed offline, and executed open loop on a
communication schedule fixed in arc length; replanning after a missed fix is
not addressed. The belief model credits every scheduled fusion at a
range-tapered weight, whereas Theurkauf et al. \cite{theurkauf2025coop}
bound the probability that a planned inter-agent measurement is
unavailable. The joint search grows exponentially in the
number of agents, and all experiments use one support. In the benchmark it
cost 1.3 to 1.9 times the wall clock of the single-agent reference on
average, most of it fixed overhead, but the exponent will bind for larger
teams. Finally, the environment is planar with static convex obstacles of
known position covariance, and the vehicle is modelled by a turn-radius
bound alone.

% ===========================================================================
\section{CONCLUSION}
% ===========================================================================
We presented a centralized planner for a primary agent that must reach a
goal in minimum time in a GPS-denied environment while support agents are
routed to keep its localization uncertainty within a terminal bound.
Optimizing over mean trajectories and propagating covariance only for
constraint evaluation reduces the belief-space problem to a joint lattice
search followed by B-spline refinement, with chance-constrained obstacle
avoidance and a curvature bound enforced in both stages. On 50 randomized
scenarios the planner solved 80\% of the scenarios at half the reference
uncertainty against 54\% for the strongest baseline, and closed-loop Monte
Carlo trials in HoloOcean showed that the certified constraints hold on the
executed trajectories. Future work includes receding-horizon replanning,
scaling the joint search to larger teams, and hardware validation.

\bibliographystyle{IEEEtran}
\bibliography{refs2}

\end{document}
